\documentclass[11pt,a4paper]{article}

\usepackage[utf8]{inputenc}
\usepackage[T1]{fontenc}
\usepackage[french]{babel}
\usepackage[a4paper,margin=2.2cm]{geometry}
\usepackage{amsmath,amssymb,amsthm,mathtools}
\usepackage{lmodern}
\usepackage{enumitem}
\usepackage{hyperref}
\usepackage{xcolor}
\usepackage{fancyhdr}
\usepackage{tikz}
\usetikzlibrary{arrows.meta,calc}

\hypersetup{
    colorlinks=true,
    linkcolor=black,
    urlcolor=blue
}

\setlength{\parindent}{0pt}
\setlength{\parskip}{0.6em}

\pagestyle{fancy}
\fancyhf{}
\fancyhead[L]{Cours --- Nombres complexes}
\fancyhead[R]{Terminale $\rightarrow$ 1re année de prépa}
\fancyfoot[C]{\thepage}

\newtheorem{definition}{Définition}[section]
\newtheorem{proposition}[definition]{Proposition}
\newtheorem{theoreme}[definition]{Théorème}
\newtheorem{corollaire}[definition]{Corollaire}
\newtheorem{remarque}[definition]{Remarque}
\newtheorem{exemple}[definition]{Exemple}
\newtheorem*{demonstration}{Démonstration}
\newtheorem*{methode}{Méthode}

\newcommand{\R}{\mathbb{R}}
\newcommand{\C}{\mathbb{C}}
\newcommand{\U}{\mathcal{U}}
\newcommand{\ii}{\mathrm{i}}
\newcommand{\e}{\mathrm{e}}
\newcommand{\dd}{\,\mathrm{d}}
\newcommand{\Ree}{\mathrm{Re}}
\newcommand{\Imm}{\mathrm{Im}}
\newcommand{\Arg}{\mathrm{Arg}}
\newcommand{\argp}{\mathrm{arg}}

\title{\textbf{Cours complet sur les nombres complexes}\\
\large De la Terminale à la première année de prépa}
\author{}
\date{}

\begin{document}

\maketitle
\tableofcontents
\bigskip

\section{Pourquoi introduire les nombres complexes ?}

Dans $\R$, tout carré est positif ou nul. En particulier, l'équation
\[
x^2+1=0
\]
n'admet aucune solution réelle. Pour pouvoir résoudre ce type d'équation, on introduit un nouveau nombre, noté $\ii$, vérifiant
\[
\ii^2=-1.
\]

\begin{definition}[Nombre complexe]
On appelle \emph{nombre complexe} tout nombre qui s'écrit sous la forme
\[
z=a+\ii b
\qquad \text{avec } a,b\in \R.
\]
L'ensemble des nombres complexes est noté $\C$.
\end{definition}

\begin{remarque}
Les nombres réels sont des nombres complexes particuliers : tout réel $a$ s'écrit
\[
a=a+\ii 0.
\]
\end{remarque}

\section{Forme algébrique et calculs}

\subsection{Partie réelle et partie imaginaire}

\begin{definition}[Écriture algébrique]
Si
\[
z=a+\ii b,
\]
on appelle :
\[
\Ree(z)=a \quad \text{la partie réelle de } z,
\qquad
\Imm(z)=b \quad \text{la partie imaginaire de } z.
\]
\end{definition}

\begin{remarque}
La \emph{partie imaginaire} est un réel. Ce n'est pas le terme en $\ii b$, mais bien le coefficient $b$.
\end{remarque}

\begin{definition}[Imaginaire pur]
On dit que $z$ est un \emph{imaginaire pur} si $\Ree(z)=0$, c'est-à-dire si
\[
z=\ii b \quad \text{avec } b\in \R.
\]
\end{definition}

\subsection{Unicité de l'écriture algébrique}

\begin{proposition}[Unicité]
Tout nombre complexe possède une unique écriture sous la forme
\[
z=a+\ii b
\quad \text{avec } a,b\in \R.
\]
\end{proposition}

\begin{demonstration}
Supposons
\[
a+\ii b=\alpha+\ii \beta
\qquad \text{avec } a,b,\alpha,\beta\in \R.
\]
Alors
\[
(a-\alpha)+\ii(b-\beta)=0.
\]
Donc
\[
a-\alpha=-\ii(b-\beta).
\]
Si $b\neq \beta$, on obtiendrait
\[
\ii=-\frac{a-\alpha}{b-\beta}\in \R,
\]
ce qui est impossible puisque $\ii^2=-1$ et aucun réel ne vérifie cela.

Donc $b=\beta$, puis nécessairement $a=\alpha$.

Ainsi l'écriture est unique.
\end{demonstration}

\subsection{Somme, produit, inverse}

Soient
\[
z_1=a+\ii b,\qquad z_2=c+\ii d.
\]
Alors
\[
z_1+z_2=(a+c)+\ii(b+d),
\]
et
\[
z_1z_2=(a+\ii b)(c+\ii d)=(ac-bd)+\ii(ad+bc).
\]

\begin{proposition}[Formules sur partie réelle et imaginaire]
Pour tous $z_1,z_2\in \C$,
\[
\Ree(z_1+z_2)=\Ree(z_1)+\Ree(z_2),\qquad
\Imm(z_1+z_2)=\Imm(z_1)+\Imm(z_2),
\]
et
\[
\Ree(z_1z_2)=\Ree(z_1)\Ree(z_2)-\Imm(z_1)\Imm(z_2),
\]
\[
\Imm(z_1z_2)=\Ree(z_1)\Imm(z_2)+\Imm(z_1)\Ree(z_2).
\]
\end{proposition}

\begin{demonstration}
Écrivons
\[
z_1=a+\ii b,\qquad z_2=c+\ii d.
\]
Alors
\[
z_1+z_2=(a+c)+\ii(b+d),
\]
d'où les deux premières formules.

De plus,
\[
z_1z_2=(a+\ii b)(c+\ii d)=ac+\ii ad+\ii bc+\ii^2bd.
\]
Comme $\ii^2=-1$, on obtient
\[
z_1z_2=(ac-bd)+\ii(ad+bc),
\]
ce qui donne les deux autres relations.
\end{demonstration}

\begin{proposition}[Inverse d'un complexe non nul]
Si
\[
z=a+\ii b\neq 0,
\]
alors
\[
\frac{1}{z}=\frac{a-\ii b}{a^2+b^2}.
\]
\end{proposition}

\begin{demonstration}
On calcule
\[
(a+\ii b)(a-\ii b)=a^2-(\ii b)^2=a^2+b^2.
\]
Comme $z\neq 0$, on n'a pas simultanément $a=0$ et $b=0$, donc
\[
a^2+b^2>0.
\]
Ainsi,
\[
\frac{1}{a+\ii b}
=\frac{a-\ii b}{(a+\ii b)(a-\ii b)}
=\frac{a-\ii b}{a^2+b^2}.
\]
\end{demonstration}

\begin{methode}
Pour mettre un quotient sous forme algébrique, on multiplie numérateur et dénominateur par le \emph{conjugué du dénominateur}.
\end{methode}

\begin{exemple}
\[
\frac{1}{2+3\ii}
=\frac{2-3\ii}{(2+3\ii)(2-3\ii)}
=\frac{2-3\ii}{13}.
\]
\end{exemple}

\section{Interprétation géométrique : affixe d'un point}

On se place dans un repère orthonormé $(O,\vec{u},\vec{v})$.

\begin{definition}[Affixe]
Au point $M(a,b)$ du plan, on associe le nombre complexe
\[
z=a+\ii b.
\]
On dit que $z$ est l'\emph{affixe} du point $M$.
\end{definition}

Ainsi :
\[
\Ree(z)=\text{abscisse de }M,
\qquad
\Imm(z)=\text{ordonnée de }M.
\]

\begin{proposition}[Affixe d'un vecteur]
Si $A$ et $B$ ont pour affixes respectives $z_A$ et $z_B$, alors le vecteur $\overrightarrow{AB}$ a pour affixe
\[
z_B-z_A.
\]
\end{proposition}

\begin{demonstration}
Si
\[
z_A=x_A+\ii y_A,\qquad z_B=x_B+\ii y_B,
\]
alors
\[
\overrightarrow{AB}=(x_B-x_A,\;y_B-y_A),
\]
donc son affixe est
\[
(x_B-x_A)+\ii(y_B-y_A)=z_B-z_A.
\]
\end{demonstration}

\begin{proposition}[Milieu]
Le milieu du segment $[AB]$ a pour affixe
\[
\frac{z_A+z_B}{2}.
\]
\end{proposition}

\begin{demonstration}
Le milieu $I$ a pour coordonnées
\[
\left(\frac{x_A+x_B}{2},\frac{y_A+y_B}{2}\right),
\]
donc pour affixe
\[
\frac{x_A+x_B}{2}+\ii\frac{y_A+y_B}{2}
=\frac{z_A+z_B}{2}.
\]
\end{demonstration}

\section{Conjugué d'un nombre complexe}

\begin{definition}[Conjugué]
Si
\[
z=a+\ii b,
\]
on appelle \emph{conjugué} de $z$ le nombre
\[
\overline{z}=a-\ii b.
\]
\end{definition}

\begin{remarque}
Géométriquement, le point d'affixe $\overline{z}$ est le symétrique du point d'affixe $z$ par rapport à l'axe réel.
\end{remarque}

\begin{proposition}
Pour tout $z\in \C$,
\[
z\in \R \iff z=\overline{z},
\qquad
z\in \ii\R \iff z=-\overline{z}.
\]
\end{proposition}

\begin{demonstration}
Si $z=a+\ii b$, alors
\[
\overline{z}=a-\ii b.
\]
Donc
\[
z=\overline{z}
\iff a+\ii b=a-\ii b
\iff 2\ii b=0
\iff b=0,
\]
c'est-à-dire $z\in \R$.

De même,
\[
z=-\overline{z}
\iff a+\ii b=-a+\ii b
\iff 2a=0
\iff a=0,
\]
donc $z\in \ii\R$.
\end{demonstration}

\begin{proposition}[Opérations et conjugué]
Pour tous $z,z'\in \C$,
\[
\overline{z+z'}=\overline{z}+\overline{z'},
\qquad
\overline{zz'}=\overline{z}\,\overline{z'}.
\]
Si $z'\neq 0$, alors
\[
\overline{\frac{z}{z'}}=\frac{\overline{z}}{\overline{z'}}.
\]
\end{proposition}

\begin{demonstration}
Écrivons
\[
z=a+\ii b,\qquad z'=c+\ii d.
\]
Alors
\[
\overline{z+z'}=\overline{(a+c)+\ii(b+d)}=(a+c)-\ii(b+d)
=\overline{z}+\overline{z'}.
\]
Pour le produit,
\[
zz'=(ac-bd)+\ii(ad+bc),
\]
donc
\[
\overline{zz'}=(ac-bd)-\ii(ad+bc).
\]
D'autre part,
\[
\overline{z}\,\overline{z'}=(a-\ii b)(c-\ii d)=(ac-bd)-\ii(ad+bc).
\]
Les deux sont donc égaux.

Enfin, si $z'\neq 0$,
\[
\overline{\frac{z}{z'}}=\overline{z\cdot \frac1{z'}}=\overline{z}\cdot \overline{\frac1{z'}}
=\overline{z}\cdot \frac1{\overline{z'}}
=\frac{\overline{z}}{\overline{z'}}.
\]
\end{demonstration}

\begin{proposition}[Lien avec partie réelle et partie imaginaire]
Pour tout $z\in \C$,
\[
\Ree(z)=\frac{z+\overline{z}}{2},
\qquad
\Imm(z)=\frac{z-\overline{z}}{2\ii}.
\]
\end{proposition}

\begin{demonstration}
Si $z=a+\ii b$, alors
\[
z+\overline{z}=(a+\ii b)+(a-\ii b)=2a,
\]
donc
\[
\frac{z+\overline{z}}{2}=a=\Ree(z).
\]
De même,
\[
z-\overline{z}=(a+\ii b)-(a-\ii b)=2\ii b,
\]
donc
\[
\frac{z-\overline{z}}{2\ii}=b=\Imm(z).
\]
\end{demonstration}

\section{Module d'un nombre complexe}

\begin{definition}[Module]
Si
\[
z=a+\ii b,
\]
on appelle \emph{module} de $z$ le réel positif
\[
|z|=\sqrt{a^2+b^2}.
\]
\end{definition}

\begin{proposition}[Relation fondamentale]
Pour tout $z\in \C$,
\[
z\overline{z}=|z|^2.
\]
\end{proposition}

\begin{demonstration}
Si $z=a+\ii b$, alors
\[
z\overline{z}=(a+\ii b)(a-\ii b)=a^2+b^2=|z|^2.
\]
\end{demonstration}

\begin{corollaire}
Si $z\neq 0$, alors
\[
\frac1z=\frac{\overline{z}}{|z|^2}.
\]
\end{corollaire}

\begin{demonstration}
Comme
\[
z\overline{z}=|z|^2,
\]
on a
\[
z\cdot \frac{\overline{z}}{|z|^2}=1.
\]
Donc $\overline{z}/|z|^2$ est l'inverse de $z$.
\end{demonstration}

\begin{proposition}[Premières propriétés]
Pour tout $z\in \C$,
\[
|z|\geq 0,
\qquad
|z|=0 \iff z=0,
\qquad
|\overline{z}|=|z|.
\]
\end{proposition}

\begin{demonstration}
Par définition,
\[
|z|=\sqrt{a^2+b^2}\geq 0.
\]
De plus,
\[
|z|=0 \iff a^2+b^2=0.
\]
Comme $a^2\geq 0$ et $b^2\geq 0$, cela équivaut à
\[
a=0 \quad \text{et} \quad b=0,
\]
donc $z=0$.

Enfin,
\[
|\overline{z}|=\sqrt{a^2+(-b)^2}=\sqrt{a^2+b^2}=|z|.
\]
\end{demonstration}

\begin{proposition}[Module d'un produit et d'un quotient]
Pour tous $z,z'\in \C$,
\[
|zz'|=|z||z'|.
\]
Si de plus $z'\neq 0$,
\[
\left|\frac{z}{z'}\right|=\frac{|z|}{|z'|}.
\]
\end{proposition}

\begin{demonstration}
On utilise la relation fondamentale :
\[
|zz'|^2=(zz')\overline{(zz')}=zz'\,\overline{z'}\,\overline{z}
=(z\overline{z})(z'\overline{z'})=|z|^2|z'|^2.
\]
Comme les modules sont positifs,
\[
|zz'|=|z||z'|.
\]
Puis, si $z'\neq 0$,
\[
\left|\frac{z}{z'}\right|
=\left|z\cdot \frac1{z'}\right|
=|z|\cdot \left|\frac1{z'}\right|
=|z|\cdot \frac1{|z'|}
=\frac{|z|}{|z'|}.
\]
\end{demonstration}

\begin{proposition}[Comparaison avec la partie réelle et la partie imaginaire]
Pour tout $z\in \C$,
\[
|\Ree(z)|\leq |z|,
\qquad
|\Imm(z)|\leq |z|.
\]
\end{proposition}

\begin{demonstration}
Si $z=a+\ii b$, alors
\[
a^2\leq a^2+b^2,
\qquad
b^2\leq a^2+b^2.
\]
En prenant la racine carrée,
\[
|a|\leq \sqrt{a^2+b^2}=|z|,
\qquad
|b|\leq \sqrt{a^2+b^2}=|z|.
\]
\end{demonstration}

\subsection{Interprétation géométrique}

Si $M$ est le point d'affixe $z=a+\ii b$, alors
\[
|z|=\sqrt{a^2+b^2}=OM.
\]

\begin{proposition}[Distance entre deux points]
Si $A$ et $B$ ont pour affixes $z_A$ et $z_B$, alors
\[
AB=|z_B-z_A|.
\]
\end{proposition}

\begin{demonstration}
Le vecteur $\overrightarrow{AB}$ a pour affixe $z_B-z_A$, donc ses coordonnées sont
\[
(\Ree(z_B-z_A),\, \Imm(z_B-z_A)).
\]
Sa norme vaut alors
\[
\sqrt{(\Ree(z_B-z_A))^2+(\Imm(z_B-z_A))^2}=|z_B-z_A|.
\]
\end{demonstration}

\section{Inégalité triangulaire}

\begin{proposition}[Inégalité triangulaire]
Pour tous $z,z'\in \C$,
\[
|z+z'|\leq |z|+|z'|.
\]
\end{proposition}

\begin{demonstration}
On calcule :
\[
|z+z'|^2=(z+z')(\overline{z}+\overline{z'})
=|z|^2+|z'|^2+z\overline{z'}+\overline{z}z'.
\]
Or
\[
z\overline{z'}+\overline{z}z'=2\Ree(z\overline{z'}).
\]
Donc
\[
|z+z'|^2=|z|^2+|z'|^2+2\Ree(z\overline{z'}).
\]
Comme
\[
\Ree(w)\leq |w| \quad \text{pour tout } w\in \C,
\]
on obtient
\[
|z+z'|^2\leq |z|^2+|z'|^2+2|z\overline{z'}|.
\]
Puis
\[
|z\overline{z'}|=|z||\overline{z'}|=|z||z'|.
\]
Ainsi
\[
|z+z'|^2\leq |z|^2+|z'|^2+2|z||z'|=(|z|+|z'|)^2.
\]
Comme les deux membres sont positifs,
\[
|z+z'|\leq |z|+|z'|.
\]
\end{demonstration}

\begin{corollaire}[Deuxième inégalité triangulaire]
Pour tous $z,z'\in \C$,
\[
\bigl||z|-|z'|\bigr|\leq |z-z'|.
\]
\end{corollaire}

\begin{demonstration}
On écrit
\[
z=(z-z')+z'.
\]
Donc, par l'inégalité triangulaire,
\[
|z|\leq |z-z'|+|z'|,
\]
d'où
\[
|z|-|z'|\leq |z-z'|.
\]
En échangeant $z$ et $z'$, on obtient
\[
|z'|-|z|\leq |z-z'|.
\]
Les deux inégalités donnent
\[
\bigl||z|-|z'|\bigr|\leq |z-z'|.
\]
\end{demonstration}

\section{Complexes de module 1 et écriture exponentielle}

\subsection{Le cercle unité}

\begin{definition}
On note
\[
\U=\{z\in \C \mid |z|=1\}.
\]
C'est l'ensemble des complexes de module $1$, c'est-à-dire le cercle trigonométrique.
\end{definition}

\begin{definition}[Exponentielle imaginaire]
Pour tout réel $\theta$, on pose
\[
\e^{\ii\theta}=\cos\theta+\ii\sin\theta.
\]
\end{definition}

\begin{proposition}
Pour tout réel $\theta$,
\[
\left|\e^{\ii\theta}\right|=1.
\]
\end{proposition}

\begin{demonstration}
\[
\left|\e^{\ii\theta}\right|
=|\cos\theta+\ii\sin\theta|
=\sqrt{\cos^2\theta+\sin^2\theta}
=1.
\]
\end{demonstration}

Donc tous les nombres $\e^{\ii\theta}$ sont sur le cercle unité.

\subsection{Produit des exponentielles}

\begin{proposition}
Pour tous réels $a,b$,
\[
\e^{\ii(a+b)}=\e^{\ii a}\e^{\ii b}.
\]
\end{proposition}

\begin{demonstration}
\[
\e^{\ii a}\e^{\ii b}
=(\cos a+\ii\sin a)(\cos b+\ii\sin b).
\]
En développant :
\[
\e^{\ii a}\e^{\ii b}
=(\cos a\cos b-\sin a\sin b)+\ii(\sin a\cos b+\cos a\sin b).
\]
Or les formules d'addition donnent :
\[
\cos(a+b)=\cos a\cos b-\sin a\sin b,
\]
\[
\sin(a+b)=\sin a\cos b+\cos a\sin b.
\]
Donc
\[
\e^{\ii a}\e^{\ii b}
=\cos(a+b)+\ii\sin(a+b)
=\e^{\ii(a+b)}.
\]
\end{demonstration}

\begin{corollaire}
Pour tout réel $a$,
\[
\e^{-\ii a}=\frac{1}{\e^{\ii a}}=\overline{\e^{\ii a}}.
\]
\end{corollaire}

\begin{demonstration}
En appliquant la proposition précédente à $a$ et $-a$,
\[
\e^{\ii a}\e^{-\ii a}=\e^0=1.
\]
Donc
\[
\e^{-\ii a}=\frac1{\e^{\ii a}}.
\]
Par ailleurs,
\[
\overline{\e^{\ii a}}=\cos a-\ii\sin a=\cos(-a)+\ii\sin(-a)=\e^{-\ii a}.
\]
\end{demonstration}

\section{Formules d'Euler et formule de Moivre}

\begin{proposition}[Formules d'Euler]
Pour tout réel $\theta$,
\[
\cos\theta=\frac{\e^{\ii\theta}+\e^{-\ii\theta}}{2},
\qquad
\sin\theta=\frac{\e^{\ii\theta}-\e^{-\ii\theta}}{2\ii}.
\]
\end{proposition}

\begin{demonstration}
On a
\[
\e^{\ii\theta}=\cos\theta+\ii\sin\theta
\qquad \text{et} \qquad
\e^{-\ii\theta}=\cos\theta-\ii\sin\theta.
\]
En additionnant,
\[
\e^{\ii\theta}+\e^{-\ii\theta}=2\cos\theta.
\]
En soustrayant,
\[
\e^{\ii\theta}-\e^{-\ii\theta}=2\ii\sin\theta.
\]
D'où les formules.
\end{demonstration}

\begin{theoreme}[Formule de Moivre]
Pour tout entier $n\in \mathbb{Z}$ et tout réel $\theta$,
\[
(\cos\theta+\ii\sin\theta)^n=\cos(n\theta)+\ii\sin(n\theta).
\]
\end{theoreme}

\begin{demonstration}
Comme
\[
\cos\theta+\ii\sin\theta=\e^{\ii\theta},
\]
on a
\[
(\cos\theta+\ii\sin\theta)^n=(\e^{\ii\theta})^n=\e^{\ii n\theta}
=\cos(n\theta)+\ii\sin(n\theta).
\]
\end{demonstration}

\begin{exemple}[Linéarisation]
\[
\sin^3 x
=\left(\frac{\e^{\ii x}-\e^{-\ii x}}{2\ii}\right)^3
=\frac{3\sin x-\sin(3x)}{4}.
\]
\end{exemple}

\section{Forme trigonométrique et forme exponentielle}

\begin{theoreme}[Existence]
Tout complexe non nul $z$ s'écrit
\[
z=r(\cos\theta+\ii\sin\theta)=r\e^{\ii\theta},
\]
avec $r>0$ et $\theta\in \R$.
\end{theoreme}

\begin{demonstration}
Soit
\[
z=a+\ii b\neq 0.
\]
On pose
\[
r=|z|=\sqrt{a^2+b^2}>0.
\]
Alors
\[
\left(\frac{a}{r}\right)^2+\left(\frac{b}{r}\right)^2=1.
\]
Le point de coordonnées
\[
\left(\frac{a}{r},\frac{b}{r}\right)
\]
appartient donc au cercle trigonométrique. Il existe alors un réel $\theta$ tel que
\[
\cos\theta=\frac{a}{r},
\qquad
\sin\theta=\frac{b}{r}.
\]
Donc
\[
a=r\cos\theta,\qquad b=r\sin\theta,
\]
et ainsi
\[
z=a+\ii b=r(\cos\theta+\ii\sin\theta)=r\e^{\ii\theta}.
\]
\end{demonstration}

\begin{definition}[Argument]
Si
\[
z=|z|\e^{\ii\theta},
\]
on dit que $\theta$ est un \emph{argument} de $z$.
\end{definition}

\begin{proposition}[Non-unicité]
Si $\theta$ est un argument de $z\neq 0$, alors les arguments de $z$ sont exactement les
\[
\theta+2k\pi,\qquad k\in \mathbb{Z}.
\]
\end{proposition}

\begin{demonstration}
Comme
\[
\e^{\ii(\theta+2k\pi)}=\e^{\ii\theta}\e^{2k\ii\pi}=\e^{\ii\theta},
\]
tout $\theta+2k\pi$ est aussi un argument.

Réciproquement, si $\theta'$ est un autre argument, alors
\[
|z|\e^{\ii\theta}=|z|\e^{\ii\theta'}.
\]
Comme $|z|>0$,
\[
\e^{\ii\theta}=\e^{\ii\theta'}.
\]
Donc
\[
\cos\theta=\cos\theta',
\qquad
\sin\theta=\sin\theta',
\]
ce qui implique
\[
\theta'-\theta\in 2\pi\mathbb{Z}.
\]
\end{demonstration}

\begin{remarque}
On peut choisir un \emph{argument principal} dans $]-\pi,\pi]$.
\end{remarque}

\section{Calculs en forme exponentielle}

Soient
\[
z=r\e^{\ii\theta},
\qquad
z'=r'\e^{\ii\theta'}.
\]

Alors
\[
zz'=rr'\e^{\ii(\theta+\theta')},
\qquad
\frac{z}{z'}=\frac{r}{r'}\e^{\ii(\theta-\theta')}\quad (z'\neq 0).
\]

\begin{proposition}[Arguments]
Pour $z,z'\neq 0$,
\[
\argp(zz')\equiv \argp(z)+\argp(z')\ [2\pi],
\]
\[
\argp\left(\frac{z}{z'}\right)\equiv \argp(z)-\argp(z')\ [2\pi].
\]
\end{proposition}

\begin{demonstration}
Si
\[
z=|z|\e^{\ii\theta},
\qquad
z'=|z'|\e^{\ii\theta'},
\]
alors
\[
zz'=|z||z'|\e^{\ii(\theta+\theta')},
\]
donc $\theta+\theta'$ est un argument de $zz'$.

De même,
\[
\frac{z}{z'}=\frac{|z|}{|z'|}\e^{\ii(\theta-\theta')},
\]
donc $\theta-\theta'$ est un argument de $z/z'$.
\end{demonstration}

\section{Transformer $a\cos t+b\sin t$}

\begin{proposition}
Soient $a,b\in \R$, non tous deux nuls. Il existe $R>0$ et $\varphi\in \R$ tels que
\[
a\cos t+b\sin t=R\cos(t-\varphi)
\quad \text{pour tout } t\in \R.
\]
On peut prendre
\[
R=\sqrt{a^2+b^2}.
\]
\end{proposition}

\begin{demonstration}
Posons
\[
R=\sqrt{a^2+b^2}>0.
\]
Alors
\[
\left(\frac{a}{R}\right)^2+\left(\frac{b}{R}\right)^2=1.
\]
Il existe donc $\varphi\in \R$ tel que
\[
\cos\varphi=\frac{a}{R},
\qquad
\sin\varphi=\frac{b}{R}.
\]
Ainsi
\[
a=R\cos\varphi,\qquad b=R\sin\varphi.
\]
Donc
\[
a\cos t+b\sin t
=R(\cos\varphi\cos t+\sin\varphi\sin t).
\]
Or
\[
\cos(t-\varphi)=\cos t\cos\varphi+\sin t\sin\varphi.
\]
Donc
\[
a\cos t+b\sin t=R\cos(t-\varphi).
\]
\end{demonstration}

\section{Équations du second degré et racines carrées}

\subsection{Racines carrées d'un réel négatif}

Si $a<0$, alors
\[
(\ii\sqrt{-a})^2=a,
\qquad
(-\ii\sqrt{-a})^2=a.
\]
Donc tout réel négatif possède deux racines carrées complexes.

\subsection{Racines carrées d'un complexe}

\begin{proposition}
Tout complexe non nul $Z=\rho \e^{\ii\theta}$ admet exactement deux racines carrées :
\[
\sqrt{\rho}\,\e^{\ii\theta/2}
\qquad \text{et} \qquad
-\sqrt{\rho}\,\e^{\ii\theta/2}.
\]
\end{proposition}

\begin{demonstration}
Cherchons $z$ sous la forme
\[
z=r\e^{\ii\varphi}
\]
avec $r>0$.
L'équation
\[
z^2=Z
\]
devient
\[
r^2\e^{2\ii\varphi}=\rho \e^{\ii\theta}.
\]
On doit donc avoir
\[
r^2=\rho
\qquad \text{et} \qquad
2\varphi\equiv \theta \ [2\pi].
\]
La première condition donne
\[
r=\sqrt{\rho}.
\]
La seconde donne
\[
\varphi\equiv \frac{\theta}{2} \ [\pi].
\]
Il y a donc exactement deux solutions opposées :
\[
\sqrt{\rho}\,\e^{\ii\theta/2}
\quad \text{et} \quad
-\sqrt{\rho}\,\e^{\ii\theta/2}.
\]
\end{demonstration}

\subsection{Équation du second degré}

\begin{theoreme}
Soient $a,b,c\in \C$ avec $a\neq 0$. L'équation
\[
az^2+bz+c=0
\]
a pour discriminant
\[
\Delta=b^2-4ac.
\]
Si $\delta$ est une racine carrée de $\Delta$, alors les solutions sont
\[
z=\frac{-b\pm \delta}{2a}.
\]
\end{theoreme}

\begin{demonstration}
On part de
\[
az^2+bz+c=0.
\]
Comme $a\neq 0$,
\[
z^2+\frac{b}{a}z+\frac{c}{a}=0.
\]
On complète le carré :
\[
z^2+\frac{b}{a}z
=
\left(z+\frac{b}{2a}\right)^2-\frac{b^2}{4a^2}.
\]
Donc
\[
\left(z+\frac{b}{2a}\right)^2-\frac{b^2}{4a^2}+\frac{c}{a}=0,
\]
soit
\[
\left(z+\frac{b}{2a}\right)^2
=\frac{b^2-4ac}{4a^2}
=\frac{\Delta}{4a^2}.
\]
Si $\delta^2=\Delta$, alors
\[
z+\frac{b}{2a}=\pm \frac{\delta}{2a},
\]
d'où
\[
z=\frac{-b\pm \delta}{2a}.
\]
\end{demonstration}

\section{Racines $n$-ièmes de l'unité}

\begin{definition}
On appelle \emph{racine $n$-ième de l'unité} tout complexe $z$ tel que
\[
z^n=1.
\]
\end{definition}

\begin{theoreme}
Les racines $n$-ièmes de l'unité sont exactement les
\[
\omega_k=\e^{2k\ii\pi/n},
\qquad k=0,1,\dots,n-1.
\]
\end{theoreme}

\begin{demonstration}
Si $z^n=1$, alors $z\neq 0$, donc
\[
z=r\e^{\ii\theta}.
\]
Alors
\[
z^n=r^n\e^{\ii n\theta}=1=\e^{2m\ii\pi}
\]
pour un certain entier $m$.

On doit donc avoir
\[
r^n=1 \quad \Rightarrow \quad r=1,
\]
et
\[
n\theta\equiv 0 \ [2\pi].
\]
Ainsi
\[
\theta=\frac{2k\pi}{n}
\]
pour un entier $k$.

Donc toute racine s'écrit
\[
z=\e^{2k\ii\pi/n}.
\]
Comme ajouter $n$ à $k$ ne change pas le nombre, on obtient exactement
\[
k=0,1,\dots,n-1.
\]
\end{demonstration}

\begin{remarque}
Ces racines sont les sommets d'un polygone régulier à $n$ côtés inscrit dans le cercle unité.
\end{remarque}

\section{Racines $n$-ièmes d'un complexe non nul}

\begin{theoreme}
Soit
\[
Z=\rho \e^{\ii\theta}
\qquad (\rho>0).
\]
Les solutions de
\[
z^n=Z
\]
sont
\[
z_k=\sqrt[n]{\rho}\;\e^{\ii(\theta+2k\pi)/n},
\qquad k=0,1,\dots,n-1.
\]
\end{theoreme}

\begin{demonstration}
On cherche
\[
z=r\e^{\ii\varphi}.
\]
Alors
\[
z^n=r^n\e^{\ii n\varphi}=\rho \e^{\ii\theta}.
\]
Donc
\[
r^n=\rho
\qquad \text{et} \qquad
n\varphi\equiv \theta \ [2\pi].
\]
Ainsi
\[
r=\sqrt[n]{\rho},
\qquad
\varphi=\frac{\theta+2k\pi}{n}.
\]
On obtient donc
\[
z_k=\sqrt[n]{\rho}\;\e^{\ii(\theta+2k\pi)/n}.
\]
Il y en a exactement $n$, pour $k=0,\dots,n-1$.
\end{demonstration}

\section{Bilan de méthode}

\subsection*{En Terminale}
On maîtrise surtout :
\begin{itemize}[leftmargin=1.2cm]
    \item la forme algébrique $a+\ii b$ ;
    \item les calculs de somme, produit, quotient ;
    \item le conjugué ;
    \item le module ;
    \item la représentation géométrique ;
    \item l'écriture trigonométrique ;
    \item les racines carrées et quelques équations simples.
\end{itemize}

\subsection*{En première année de prépa}
On approfondit :
\begin{itemize}[leftmargin=1.2cm]
    \item l'usage systématique de la forme exponentielle ;
    \item les arguments ;
    \item les formules d'Euler et de Moivre ;
    \item la linéarisation/trigonométrie ;
    \item les racines $n$-ièmes de l'unité ;
    \item la résolution générale de $z^n=Z$ ;
    \item les applications géométriques.
\end{itemize}

\section{Exercices d'application}

\begin{exemple}
Mettre sous forme algébrique :
\[
\frac{3+2\ii}{1-\ii}.
\]
\end{exemple}

\begin{exemple}
Déterminer le module et un argument de
\[
z=-1+\ii\sqrt{3}.
\]
\end{exemple}

\begin{exemple}
Résoudre dans $\C$ :
\[
z^2-(1+4\ii)z+(3\ii-3)=0.
\]
\end{exemple}

\begin{exemple}
Déterminer les racines cubiques de $1$ puis les racines quatrièmes de $-16$.
\end{exemple}

\bigskip
\begin{center}
\textbf{Fin du cours}
\end{center}

\end{document}